% Part: computability
% Chapter: computability-theory
% Section: reducibility

\documentclass[../../../include/open-logic-section]{subfiles}

\begin{document}

\olfileid{cmp}{thy}{red}
\olsection{کاهش‌پذیری}

\begin{explain}
اکنون می‌دانیم دست‌کم یک مجموعه، یعنی $K_0$، وجود دارد که
محاسبه‌پذیرِ برشمردنی است اما محاسبه‌پذیر نیست. باید روشن باشد که
مجموعه‌های دیگری نیز چنین‌اند. روش کاهش‌پذیری ابزاری نیرومند برای
نشان‌دادن این ویژگی‌ها در مجموعه‌های دیگر فراهم می‌کند، بی‌آنکه ناچار
باشیم پیوسته به اصول نخستین بازگردیم.

به‌طور کلی، «کاهش» مجموعهٔ $A$ به مجموعهٔ~$B$ روشی است برای تبدیل
پاسخ پرسش‌های عضویت !!{element}s در~$B$ به پاسخ پرسش‌های عضویت
!!{element}s در~$A$. ما بر مفهومی به نام «کاهش‌پذیری چندبه‌یک» تمرکز
می‌کنیم، اما مفهوم‌های بسیار دیگری از کاهش‌پذیری با ویژگی‌های گوناگون
وجود دارند. مفهوم‌های کاهش‌پذیری در مطالعهٔ پیچیدگی محاسباتی نیز
مرکزی‌اند؛ جایی که باید کارایی را هم در نظر گرفت. برای نمونه،
مجموعه‌ای «NP-کامل» نامیده می‌شود اگر در NP باشد و هر مسئلهٔ NP را
بتوان با مفهومی از کاهش همانند مفهوم زیر به آن کاهش داد، با این شرط
اضافی که کاهش در زمان چندجمله‌ای محاسبه‌پذیر باشد.

از پیش مفهوم کاهش را به‌طور ضمنی به کار برده‌ایم. مجموعهٔ~$K$ را
چنین تعریف کنید:
\[
K = \Setabs{x}{\cfind{x}(x) \fdefined},
\]
یعنی $K=\Setabs{x}{x\in W_x}$. برهان حل‌ناپذیری مسئلهٔ توقف
(\olref[hlt]{thm:halting-problem}) به مستقیم‌ترین صورت نشان می‌دهد که
$K$ محاسبه‌پذیر نیست. به یاد آورید که $K_0$ مجموعهٔ
\[
K_0 = \Setabs{\tuple{e, x}}{\cfind{e}(x) \fdefined },
\]
یعنی $K_0=\Setabs{\tuple{e,x}}{x\in W_e}$ است. هر برهان
محاسبه‌ناپذیریِ~$K$ را به‌آسانی می‌توان به برهانی برای
محاسبه‌ناپذیریِ~$K_0$ گسترش داد: اگر $K_0$ محاسبه‌پذیر بود، می‌توانستیم
صرفاً با پرسیدن اینکه آیا زوج $\tuple{x,x}$ در $K_0$ هست یا نه،
دربارهٔ عضویت !!a{element}~$x$ در $K$ تصمیم بگیریم. تابع~$f$ که $x$
را به $\tuple{x,x}$ می‌فرستد نمونه‌ای از \emph{کاهش} $K$ به $K_0$
است.
\end{explain}

\begin{defn}
بگذارید $A$ و~$B$ مجموعه‌هایی از اعداد طبیعی باشند. تابع محاسبه‌پذیر
$f\colon\Nat\to\Nat$ \emph{کاهش چندبه‌یکِ} $A$ به~$B$ است اگر و تنها
اگر به ازای هر عدد طبیعی~$x$،
\[
x \in A \quad \text{اگر و تنها اگر} \quad f(x) \in B.
\]
اگر چنین کاهش~$f$ای وجود داشته باشد، می‌گوییم $A$ به~$B$
\emph{کاهش‌پذیر چندبه‌یک} است و می‌نویسیم $A\leq_m B$. اگر $A$ به
$B$ و $B$ به~$A$ کاهش‌پذیر چندبه‌یک باشند، آن‌ها را \emph{هم‌ارز
چندبه‌یک} می‌نامیم و می‌نویسیم $A\equiv_m B$.
\end{defn}

اگر تابع $f$ در تعریف بالا اتفاقاً یک‌به‌یک باشد، می‌گوییم $A$ به~$B$
\emph{کاهش‌پذیر یک‌به‌یک} است. بیشتر کاهش‌هایی که در ادامه توصیف
می‌شوند این شرط قوی‌تر را دارند، اما از این واقعیت استفاده نخواهیم کرد.

\begin{digress}
درست است، اما به‌هیچ‌وجه بدیهی نیست، که کاهش‌پذیری یک‌به‌یک واقعاً
شرطی قوی‌تر از کاهش‌پذیری چندبه‌یک است. به بیان دیگر، مجموعه‌های
نامتناهی $A$ و~$B$ای وجود دارند که $A$ به~$B$ کاهش‌پذیر چندبه‌یک
است، اما کاهش‌پذیر یک‌به‌یک نیست.
\end{digress}

\end{document}
